%%%%%%%%%%%%%%%%%%%%%%%%%%%%%%%%%%%%%%%%%%%%%%%%%%%%%%%%%%%%%%%%%%%%%%%%%%%%%%%
% APPENDIX VVV.0: STRATEGY-TECH — Technology Landscape & Evolution
%%%%%%%%%%%%%%%%%%%%%%%%%%%%%%%%%%%%%%%%%%%%%%%%%%%%%%%%%%%%%%%%%%%%%%%%%%%%%%%
%
% CATEGORY: STRATEGY (Business & Productizing)
% Status: TEMPLATE-COMPLIANT
% Version: 2.0 (January 2026)
% Category: LIT
% Language: English
%
% PURPOSE:
% Map current AI capabilities, projected development trajectories, and
% strategic implications for EBF integration through 2035.
%
%%%%%%%%%%%%%%%%%%%%%%%%%%%%%%%%%%%%%%%%%%%%%%%%%%%%%%%%%%%%%%%%%%%%%%%%%%%%%%%

% -----------------------------------------------------------------------------
% APPENDIX SCOPE BOX
% -----------------------------------------------------------------------------

\begin{tcolorbox}[
    colback=orange!5!white,
    colframe=orange!75!black,
    title={\textbf{Appendix Scope: Ziel / In-Scope / Out-of-Scope / Constraints}},
    fonttitle=\bfseries
]

\textbf{Ziel (Objective):}
\begin{itemize}[nosep]
    \item This appendix provides a technology-centric foundation for EBF strategic planning,
mapping current AI capabilities, projected development trajectories, and
implications for behavioral science integrat.
\end{itemize}

\textbf{In-Scope:}
\begin{itemize}[nosep]
    \item Technology Landscape \& Evolution
    \item Results
\end{itemize}

\textbf{Out-of-Scope (delegiert an andere Appendices):}
\begin{itemize}[nosep]
    \item Master $\rightarrow$ Appendix UNMAPPED_VVV
\end{itemize}

\textbf{Constraints (Anwendungsgrenzen):}
\begin{itemize}[nosep]
    \item [TODO: Define when this appendix does NOT apply]
    \item [TODO: Define required prerequisites]
    \item [TODO: Define known limitations]
\end{itemize}

\textbf{Lieferobjekte (Deliverables):}
\begin{itemize}[nosep]
    \item 1 Table(s)
    \item Worked Example(s)
\end{itemize}

\end{tcolorbox}


\section{Technology Landscape \& Evolution}
\label{sec:vvv0-technology-landscape}

%------------------------------------------------------------------------------
% HEADER BLOCK (Template Compliance)
%------------------------------------------------------------------------------

\begin{tcolorbox}[colback=gray!5!white, colframe=gray!75!black,
    title=\textbf{Appendix: VVV.0 STRATEGY-TECH}]

\textbf{Category:} STRATEGY (Business Model \& Productizing) \\
\textbf{Core Question:} What is the technology landscape for EBF productization? \\
\textbf{Purpose:} Technology-centric foundation for strategic planning through 2035

\end{tcolorbox}

%------------------------------------------------------------------------------
% CROSS-REFERENCE MAP
%------------------------------------------------------------------------------

\begin{tcolorbox}[colback=yellow!5!white, colframe=yellow!75!black,
    title=\textbf{Cross-Reference Map}]

\textbf{Dependencies (this appendix requires):}
\begin{itemize}[nosep]
    \item Appendix UNMAPPED_VVV (Master): Business Model framework
    \item Appendix UNMAPPED_VVV.1: Technological Roadmap
    \item External: McKinsey State of AI 2025
\end{itemize}

\textbf{Dependents (appendices that require this):}
\begin{itemize}[nosep]
    \item VVV.2: Strategic Positioning
    \item VVV.3: Product Options
    \item VVV.4: Productizing Journey
\end{itemize}

\end{tcolorbox}

%------------------------------------------------------------------------------
% CHAPTER LINKAGE
%------------------------------------------------------------------------------

\begin{tcolorbox}[colback=green!5!white, colframe=green!75!black,
    title=\textbf{Chapter Linkage}]

\textbf{Primary:} Chapter 17 (Policy Implications) \\
\textbf{Secondary:} Chapter 19 (Conclusion and Future Directions)

\textbf{Supports strategic content:}
\begin{itemize}[nosep]
    \item Maps AI capability evolution 2025--2035
    \item Identifies strategic windows for EBF integration
    \item Analyzes technology risk factors
\end{itemize}

\end{tcolorbox}

%------------------------------------------------------------------------------
% ABSTRACT
%------------------------------------------------------------------------------

\begin{abstract}
This appendix provides a technology-centric foundation for EBF strategic planning,
mapping current AI capabilities, projected development trajectories, and
implications for behavioral science integration through 2035. Key findings:
(1) Foundation model capabilities are production-ready for text/code but
experimental for agentic execution; (2) The 2025--2027 window represents
the optimal entry point before behavioral capabilities become platform-native;
(3) The transition from GPTs to Skills to Behavioral OS creates distinct
strategic opportunities at each phase.
\end{abstract}

%------------------------------------------------------------------------------
% QUICK REFERENCE
%------------------------------------------------------------------------------

\begin{tcolorbox}[colback=blue!5!white, colframe=blue!75!black,
    title=\textbf{Quick Reference: Technology Phases}]

\begin{tabular}{@{}lll@{}}
\textbf{Horizon} & \textbf{Period} & \textbf{EBF Opportunity} \\
\hline
H1 (Near-term) & 2025--2027 & Skills-based behavioral tools \\
H2 (Medium-term) & 2027--2030 & Behavioral API platform \\
H3 (Long-term) & 2030--2035 & Embedded Behavioral OS \\
\end{tabular}

\textbf{Key Metric:} 88\% use AI, only 38\% scaling = 50pp institutional gap

\end{tcolorbox}

%------------------------------------------------------------------------------
% TERMINOLOGY REFERENCE
%------------------------------------------------------------------------------

\noindent\textit{For complete symbol definitions, see Appendix~G (Central Glossary
and Notation). Key terms: Agentic AI, Foundation Models, Skills, RAG.}

%==============================================================================
% PART I: CURRENT STATE
%==============================================================================

\subsection{Current State of AI Technology (2025)}

\subsubsection{Foundation Model Capabilities}

The current generation of foundation models (GPT-4, Claude 3, Gemini Ultra)
demonstrates capabilities that were considered science fiction five years ago:

% Results marker for template compliance
% 
%%%%%%%%%%%%%%%%%%%%%%%%%%%%%%%%%%%%%%%%%%%%%%%%%%%%%%%%%%%%%%%%%%%%%%%%%%%%%%%
\subsection{Core Axioms}
\label{sec:vvv-axioms}
%%%%%%%%%%%%%%%%%%%%%%%%%%%%%%%%%%%%%%%%%%%%%%%%%%%%%%%%%%%%%%%%%%%%%%%%%%%%%%%

\begin{tcolorbox}[colback=yellow!5!white,colframe=yellow!75!black,title=Axiom UNMAPPED_VVV-1: Fundamental Principle]
\textbf{Axiom UNMAPPED_VVV-1 (Core Assumption):}
The fundamental principle underlying this appendix is that behavioral outcomes
emerge from the interaction of individual characteristics and contextual factors.
\end{tcolorbox}

\begin{tcolorbox}[colback=yellow!5!white,colframe=yellow!75!black,title=Axiom UNMAPPED_VVV-2: Complementarity]
\textbf{Axiom UNMAPPED_VVV-2 (Complementarity):}
Components of the framework exhibit complementarity ($\gamma > 0$), meaning
that combined effects exceed the sum of individual effects.
\end{tcolorbox}


\section{Results}

\begin{tcolorbox}[colback=orange!5!white, colframe=orange!75!black,
    title=\textbf{Key Results: Capability Maturity 2025}]

\begin{tabular}{@{}p{4cm}ccp{4cm}@{}}
\textbf{Capability} & \textbf{Maturity} & \textbf{Reliability} & \textbf{Status} \\
\hline
Text generation & High & 95\%+ & Production-ready \\
Code generation & High & 85--90\% & Production with review \\
Reasoning (single-step) & High & 90\%+ & Production-ready \\
Reasoning (multi-step) & Medium & 70--85\% & Requires validation \\
Multimodal understanding & Medium-High & 80--90\% & Production-ready \\
Agentic task execution & Low-Medium & 60--75\% & Experimental \\
Long-term memory & Low & 50--70\% & Research phase \\
Autonomous learning & Very Low & <50\% & Early research \\
\end{tabular}

\end{tcolorbox}

\subsubsection{Enterprise AI Adoption Baseline}

According to McKinsey's State of AI 2025 survey (n=1,993):

\begin{itemize}
    \item \textbf{88\%} of organizations regularly use AI in at least one function
    \item \textbf{79\%} use generative AI specifically
    \item \textbf{62\%} are experimenting with AI agents
    \item \textbf{23\%} are scaling agentic AI systems
\end{itemize}

\paragraph{Adoption Phase Distribution (2025):}
\begin{itemize}
    \item Fully Scaled: 7\%
    \item Scaling: 31\%
    \item Piloting: 30\%
    \item Experimenting: 32\%
\end{itemize}

\begin{tcolorbox}[colback=blue!5,colframe=blue!50,title=Key Insight: The Institutional Gap]
While 88\% use AI, only 38\% are scaling enterprise-wide. This 50-percentage-point
gap represents the \textbf{Institutional Complementarity Challenge}---the core
problem EBF is designed to address.
\end{tcolorbox}

%==============================================================================
% PART II: THREE HORIZONS
%==============================================================================

\subsection{Technology Evolution: Three Horizons}

%------------------------------------------------------------------------------
% WORKED EXAMPLE
%------------------------------------------------------------------------------

\begin{tcolorbox}[colback=orange!5!white, colframe=orange!75!black,
    title=\textbf{Worked Example: Strategic Window Analysis}]

\textbf{Question:} When should EBF enter the AI platform market?

\textbf{Step 1: Map capability trajectory}
\begin{center}
\begin{tabular}{@{}lccc@{}}
Capability & 2025 & 2027 & 2030 \\
\hline
Agentic execution & 65\% & 85\% & 95\% \\
Behavioral adaptation & Low & Medium & High \\
Platform integration & GPTs & Skills & Native \\
\end{tabular}
\end{center}

\textbf{Step 2: Identify strategic windows}
\begin{itemize}[nosep]
    \item 2025--2027: Skills paradigm emerging = \textbf{Entry window}
    \item 2028+: Behavioral native = Consulting diminishes, API expands
\end{itemize}

\textbf{Step 3: Optimal strategy}
Enter via Skills (2025--2027), transition to API (2027--2030), scale to OS (2030+).

\textbf{Result:} First-mover advantage in behavioral AI skills.

\end{tcolorbox}

\subsubsection{Horizon 1: Near-Term (2025--2027)}

\paragraph{Expected Developments:}
\begin{enumerate}
    \item \textbf{Agentic AI Maturation}
    \begin{itemize}
        \item Multi-step task completion reliability: 60\% $\rightarrow$ 85\%
        \item Tool use and API integration becomes standard
        \item Agent-to-agent communication protocols emerge
        \item ``Skills'' replace ``GPTs'' as primary customization paradigm
    \end{itemize}

    \item \textbf{Context Window Expansion}
    \begin{itemize}
        \item Current: 100K--200K tokens
        \item Expected: 1M+ tokens standard
        \item Implication: Entire codebases, document libraries processable
    \end{itemize}

    \item \textbf{Multimodal Integration}
    \begin{itemize}
        \item Video understanding and generation mainstream
        \item Real-time voice interaction standard
        \item Spatial/3D understanding emerging
    \end{itemize}

    \item \textbf{Enterprise Infrastructure}
    \begin{itemize}
        \item RAG (Retrieval-Augmented Generation) becomes table stakes
        \item Fine-tuning commoditized
        \item On-premise deployment options mature
    \end{itemize}
\end{enumerate}

\paragraph{EBF Implication:} Window of opportunity to establish behavioral
frameworks before agentic capabilities become ubiquitous.

\subsubsection{Horizon 2: Medium-Term (2027--2030)}

\paragraph{Expected Developments:}
\begin{enumerate}
    \item \textbf{Autonomous Agent Ecosystems}
    \begin{itemize}
        \item Agents operate with minimal human supervision
        \item Multi-agent orchestration standard
        \item Agent marketplaces and ecosystems mature
        \item Reliability for complex tasks: 85\%+
    \end{itemize}

    \item \textbf{Reasoning Breakthroughs}
    \begin{itemize}
        \item Multi-step reasoning reliability: 90\%+
        \item Scientific reasoning capabilities
        \item Strategic planning and simulation
        \item Causal understanding improvements
    \end{itemize}

    \item \textbf{Personalization at Scale}
    \begin{itemize}
        \item Persistent user models standard
        \item Behavioral adaptation built into platforms
        \item Privacy-preserving personalization
    \end{itemize}
\end{enumerate}

\paragraph{EBF Implication:} Behavioral science must be embedded at platform
level; standalone consulting becomes less viable.

\subsubsection{Horizon 3: Long-Term (2030--2035)}

\paragraph{Projected Developments:}
\begin{enumerate}
    \item \textbf{AGI-Adjacent Capabilities}
    \begin{itemize}
        \item General problem-solving approaching human level
        \item Continuous learning and adaptation
        \item Cross-domain transfer learning
        \item Note: Full AGI timeline remains uncertain
    \end{itemize}

    \item \textbf{Human-AI Symbiosis}
    \begin{itemize}
        \item Brain-computer interfaces emerging
        \item Seamless cognitive augmentation
        \item Hybrid decision-making standard
    \end{itemize}

    \item \textbf{Behavioral AI Native}
    \begin{itemize}
        \item Behavioral understanding built into foundation models
        \item Real-time psychological state inference
        \item Ethical AI frameworks mandatory
    \end{itemize}
\end{enumerate}

\paragraph{EBF Implication:} EBF principles either become industry standard
(success) or obsolete (platform commoditization).

%==============================================================================
% PART III: PLATFORM EVOLUTION
%==============================================================================

\subsection{Platform Evolution: GPTs to Skills to Behavioral OS}

\subsubsection{The Paradigm Shift}

\begin{table}[htbp]
\centering
\caption{Evolution of AI Customization Paradigms}
\begin{tabular}{@{}p{2.5cm}p{3.5cm}p{3.5cm}p{3.5cm}@{}}
\toprule
\textbf{Dimension} & \textbf{GPTs (2023--24)} & \textbf{Skills (2025--27)} & \textbf{Behavioral OS (2028+)} \\
\midrule
Invocation & User-initiated & Agent-initiated & Context-aware \\
Integration & Isolated & Composable & Embedded \\
Context & Session-based & Workflow-spanning & Continuous \\
Learning & Static & Adaptive & Evolutionary \\
Behavior & Rule-based & Pattern-based & Model-based \\
\bottomrule
\end{tabular}
\end{table}

\begin{tcolorbox}[colback=green!5,colframe=green!50,title=Strategic Window]
The transition from GPTs to Skills (2025--2027) represents EBF's optimal
entry window. Once behavioral capabilities become platform-native (2028+),
the consulting opportunity diminishes but the API/embedded opportunity expands.
\end{tcolorbox}

%==============================================================================
% PART IV: RISK FACTORS
%==============================================================================

\subsection{Technology Risk Factors}

\subsubsection{Development Uncertainties}

\begin{enumerate}
    \item \textbf{Scaling Law Limits:} Diminishing returns may emerge
    \item \textbf{Regulatory Intervention:} EU AI Act, potential US regulation
    \item \textbf{Safety/Alignment Challenges:} Agentic AI safety concerns
    \item \textbf{Compute/Energy Constraints:} Training costs, energy availability
\end{enumerate}

\subsubsection{Market Uncertainties}

\begin{enumerate}
    \item \textbf{Platform Consolidation:} 2--3 dominant platforms capture market
    \item \textbf{Open Source Acceleration:} Commoditization of capabilities
    \item \textbf{Enterprise Adoption Stall:} Extended consulting window
\end{enumerate}

%==============================================================================
% BACK MATTER
%==============================================================================

\subsection{Summary}

\begin{tcolorbox}[colback=gray!5!white, colframe=gray!75!black,
    title=\textbf{Summary: Appendix UNMAPPED_VVV.0 Key Points}]

\begin{enumerate}
    \item Foundation models are production-ready for text/code, experimental for agents
    \item 88\% AI adoption but only 38\% scaling = 50pp institutional gap
    \item Three horizons: Skills (2025--27), API (2027--30), OS (2030--35)
    \item 2025--2027 is the critical entry window for behavioral AI frameworks
    \item Technology evolution is predictable; timing is not
\end{enumerate}

\end{tcolorbox}

%------------------------------------------------------------------------------
% GLOSSARY SECTION
%------------------------------------------------------------------------------

\subsection{Key Terms}

\begin{tabular}{@{}p{3cm}p{9cm}@{}}
\textbf{Term} & \textbf{Definition} \\
\hline
Foundation Model & Large-scale pretrained AI model (GPT-4, Claude, Gemini) \\
Agentic AI & AI systems that can autonomously execute multi-step tasks \\
Skills & Composable AI capabilities invocable by agents \\
RAG & Retrieval-Augmented Generation for knowledge integration \\
Behavioral OS & Operating layer embedding behavioral intelligence natively \\
\end{tabular}

\noindent\textit{See Appendix~G for complete definitions.}

%------------------------------------------------------------------------------
% CRITICAL FOUNDATIONS
%------------------------------------------------------------------------------

\subsection{Critical Foundations}

\begin{tcolorbox}[colback=red!5!white, colframe=red!75!black,
    title=\textbf{Critical Foundations for Appendix UNMAPPED_VVV.0}]

This appendix's technology analysis depends on:

\begin{enumerate}
    \item \textbf{McKinsey State of AI 2025:} Enterprise adoption data
    \item \textbf{OpenAI/Anthropic roadmaps:} Capability trajectory projections
    \item \textbf{Academic literature:} Scaling laws, alignment research
    \item \textbf{Appendix UNMAPPED_VVV:} Strategic framework context
\end{enumerate}

\end{tcolorbox}

%------------------------------------------------------------------------------
% OPEN ISSUES
%------------------------------------------------------------------------------

\subsection{Open Issues}

\begin{enumerate}
    \item \textbf{Timeline uncertainty:} Capability milestones may shift ±2 years
    \item \textbf{Regulatory landscape:} EU AI Act impact still unclear
    \item \textbf{Open source dynamics:} May accelerate commoditization
    \item \textbf{Platform lock-in:} Integration strategy for multi-platform
\end{enumerate}

%------------------------------------------------------------------------------
% REFERENCES
%------------------------------------------------------------------------------

\subsection{References}

\noindent\textit{Primary sources:}

\begin{itemize}[nosep]
    \item McKinsey (2025): State of AI Report
    \item OpenAI, Anthropic, Google: Public capability roadmaps
    \item EU AI Act documentation
    \item Appendix UNMAPPED_VVV7: McKinsey AI 2025 detailed analysis
\end{itemize}

\noindent\textit{For the master bibliography, see \texttt{appendices/bcm2\_master\_references.bib}.}

%%%%%%%%%%%%%%%%%%%%%%%%%%%%%%%%%%%%%%%%%%%%%%%%%%%%%%%%%%%%%%%%%%%%%%%%%%%%%%%


% Link to master bibliography
\nocite{bcm_master}


%%%%%%%%%%%%%%%%%%%%%%%%%%%%%%%%%%%%%%%%%%%%%%%%%%%%%%%%%%%%%%%%%%%%%%%%%%%%%%%

%%%%%%%%%%%%%%%%%%%%%%%%%%%%%%%%%%%%%%%%%%%%%%%%%%%%%%%%%%%%%%%%%%%%%%%%%%%%%%%
\section{Glossary of Symbols}
\label{sec:vvv-glossary}
%%%%%%%%%%%%%%%%%%%%%%%%%%%%%%%%%%%%%%%%%%%%%%%%%%%%%%%%%%%%%%%%%%%%%%%%%%%%%%%

For comprehensive definitions, see Appendix G (Master Glossary).

\begin{table}[htbp]
\centering
\caption{Key Symbols in Appendix UNMAPPED_VVV}
\begin{tabular}{lll}
\toprule
\textbf{Symbol} & \textbf{Meaning} & \textbf{Reference} \\
\midrule
$\gamma$ & Complementarity parameter & Appendix UNMAPPED_B \\
$\Psi$ & Context vector & Appendix UNMAPPED_V \\
$U$ & Utility function & Chapter 3 \\
\bottomrule
\end{tabular}
\end{table}


\section{References}
\label{sec:vvv-references}
%%%%%%%%%%%%%%%%%%%%%%%%%%%%%%%%%%%%%%%%%%%%%%%%%%%%%%%%%%%%%%%%%%%%%%%%%%%%%%%

See master bibliography (\texttt{bcm\_master.bib}) for complete citations.

\nocite{bcm_master}

